% ============================================================
%  preamble.tex — 日本語数学ノート テンプレート共通設定
% ============================================================

% --- 基本パッケージ ---
\usepackage{amsmath,amssymb,amsthm,mathtools}
\usepackage[dvipdfmx,hidelinks]{hyperref}
\usepackage{pxjahyper}

% --- 図・表 ---
\usepackage{tikz}
\usepackage{float}
\usetikzlibrary{decorations.pathreplacing,calc,arrows.meta,patterns,positioning}
\usepackage{pgfplots}
\pgfplotsset{compat=1.18}
\usepackage{graphicx}
\graphicspath{{fig/}}

% --- 定理環境 (tcolorbox) ---
\usepackage{tcolorbox}
\tcbuselibrary{theorems,breakable}

% Definition: 青系
\newtcbtheorem[number within=section]{definition}{Def.}{
  colback=blue!5, colframe=blue!40!black, fonttitle=\bfseries, breakable
}{def}
% Proposition: オレンジ系
\newtcbtheorem[use counter from=definition]{proposition}{Prop.}{
  colback=orange!5, colframe=orange!50!black, fonttitle=\bfseries, breakable
}{prop}
% Theorem: 赤系（章・節の目玉となる重要結果）
\newtcbtheorem[use counter from=definition]{theorem}{Thm.}{
  colback=red!5, colframe=red!40!black, fonttitle=\bfseries, breakable
}{thm}
% Claim: 緑系（Thm ほど深くはないが，証明を要する主張）
\newtcbtheorem[use counter from=definition]{claim}{Claim}{
  colback=green!5, colframe=green!45!black, fonttitle=\bfseries, breakable
}{clm}
% Lemma: シアン系（命題の証明に用いる補助結果）
\newtcbtheorem[use counter from=definition]{lemma}{Lem.}{
  colback=cyan!5, colframe=cyan!50!black, fonttitle=\bfseries, breakable
}{lem}
% Remark: 黄系
\newtcbtheorem[use counter from=definition]{remark}{Rem.}{
  colback=yellow!5, colframe=yellow!40!black, fonttitle=\bfseries, breakable
}{rem}
% Example: 灰色系
\newtcbtheorem[use counter from=definition]{example}{Ex.}{
  colback=gray!5, colframe=gray!50!black, fonttitle=\bfseries, breakable
}{ex}
% Setting: 紫系
\newtcbtheorem[use counter from=definition]{setting}{Setting}{
  colback=purple!5, colframe=purple!40!black, fonttitle=\bfseries, breakable
}{set}
% Corollary: すみれ系（定理から直ちに従う結果）
\newtcbtheorem[use counter from=definition]{corollary}{Cor.}{
  colback=violet!5, colframe=violet!40!black, fonttitle=\bfseries, breakable
}{cor}
% Memo: 薄い黒系（例より薄く、番号なし、発展的補足）
\newtcolorbox{memo*}{
  colback=black!7, colframe=black!40,
  fonttitle=\bfseries, title={Memo},
  breakable
}
% 証明環境の装飾
\tcolorboxenvironment{proof}{
  colback=white, colframe=black!20,
  leftrule=2pt, rightrule=0pt, toprule=0pt, bottomrule=0pt,
  breakable
}

% ブロックの機械可読メタデータ（proofgraph ツールが解析する）。
% PDF には一切出力されない。各ブロックの \begin{env}{title}{label} 直後に置く：
%   \blockmeta{space=polish; uses=prop:foo; route.A=prop:bar}
\newcommand{\blockmeta}[1]{}

% --- アルゴリズム環境 ---
\newcounter{algcounter}[section]
\newtcolorbox{algorithm}[1]{
  colback=white, colframe=black!70,
  fonttitle=\bfseries,
  title={Algorithm~\refstepcounter{algcounter}\thealgcounter\label{#1}},
  boxrule=0.8pt
}

% --- 便利マクロ ---
\newcommand{\R}{\mathbb{R}}
\newcommand{\Rnn}{\R_{\geq 0}} % 非負実数 R_{≥0}（コスト・結合・閉集合など 0 を含む対象）
\newcommand{\Rpos}{\R_{>0}}    % 正実数 R_{>0}（ヒストグラム・スケーリング・Gibbs 核など狭義正）
\newcommand{\N}{\mathbb{N}}
\newcommand{\E}{\mathbb{E}}
\DeclareMathOperator*{\argmin}{arg\,min}
\DeclareMathOperator*{\argmax}{arg\,max}
\DeclareMathOperator{\diag}{diag}
\DeclareMathOperator{\tr}{tr}
\DeclareMathOperator{\KL}{KL}
\DeclareMathOperator{\Id}{Id}
\newcommand{\abs}[1]{\lvert #1 \rvert}
\newcommand{\norm}[1]{\lVert #1 \rVert}
\newcommand{\inner}[2]{\langle #1,\, #2 \rangle}
\newcommand{\defeq}{\overset{\mathrm{def}}{=}}

% --- 追加の集合 ---
\newcommand{\Z}{\mathbb{Z}}
\newcommand{\Q}{\mathbb{Q}}

% --- 空間の記法 ---
\newcommand{\X}{\mathcal{X}}
\newcommand{\Y}{\mathcal{Y}}
\newcommand{\Mm}{\mathcal{M}}
\newcommand{\Bb}{\mathcal{B}}
\newcommand{\Cc}{\mathcal{C}}

% --- 確率単体 ---
\newcommand{\simplex}{\Sigma}

% --- 定数ベクトル・行列 ---
\newcommand{\ones}{\mathbf{1}}
\newcommand{\Identity}{\mathbf{I}}

% --- 測度関連 ---
\newcommand{\supp}{\operatorname{supp}}
\renewcommand{\d}{\mathrm{d}}

% --- 最適輸送固有 ---
\newcommand{\Couplings}{\mathcal{U}}
\newcommand{\CouplingsD}{\Pi}
\newcommand{\Wass}{\mathcal{W}}
\newcommand{\WassD}{\mathrm{W}}
\newcommand{\MK}{\mathcal{L}}
\newcommand{\MKD}{\mathrm{L}}
\DeclareMathOperator{\Perm}{Perm}
\DeclareMathOperator{\rank}{rank}
\newcommand{\pushforward}{_\sharp}
\newcommand{\PotentialsD}{\mathbf{R}}
\newcommand{\Potentials}{\mathcal{R}}
\newcommand{\Hb}{\mathrm{H}}
\newcommand{\KLD}{\mathrm{KL}}
% --- Ground 距離 ---
\newcommand{\dist}{d}
\newcommand{\distD}{\mathbf{D}}

% --- 範囲記法 ---
\usepackage{stmaryrd}
\newcommand{\range}[1]{\llbracket #1 \rrbracket}

% --- サイト用デモブロック指示（PDF では無効，tex2md が :::demo に変換する）---
\newcommand{\demohint}[1]{}

% --- 参考文献 ---
\usepackage[numbers,sort&compress]{natbib}
